% =====================================================================================
% Figure -- the occlusion layer beside the photoreal twin it was derived from.
% Belongs in \subsection{Capturing and Rendering the Digital Twin}, or in the occlusion
% paragraph of the system description, wherever depth occlusion is first introduced.
%
% Requires in the preamble (acmart does not load it by default):
%     \usepackage{subcaption}
%
% Expected image files, relative to the main .tex:
%     figures/occluder-mesh.jpg    (grey untextured collision/occlusion mesh)
%     figures/occluder-splat.jpg   (same site, Gaussian-splat reconstruction)
%
% ASPECT RATIO -- the two sources do not match out of the renderer:
%     occluder-mesh   1462 x 983   ~= 3:2
%     occluder-splat  1221 x 929   ~= 4:3
% Side by side at equal width they render at unequal heights. Crop both to 3:2 first;
% that costs the splat 115 px of empty foreground pavement and the mesh almost nothing:
%
%     magick figures/occluder-mesh.jpg  -gravity center -crop 1462x974+0+0 +repage \
%            figures/occluder-mesh.jpg
%     magick figures/occluder-splat.jpg -gravity north  -crop 1221x814+0+0 +repage \
%            figures/occluder-splat.jpg
%
% (-gravity north on the splat keeps the sky line and the buildings, and takes the crop
% out of the paved foreground.) After cropping, the \includegraphics calls below need no
% height or trim arguments.
%
% Reference in text as Figure~\ref{fig:occlusion}.
% A single-panel fallback is commented at the bottom of this file.
% =====================================================================================

\begin{figure*}[t]
  \centering
  \begin{subfigure}[b]{0.48\textwidth}
    \includegraphics[width=\linewidth]{figures/occluder-mesh}
    \caption{Occlusion layer.}
    \label{fig:occlusion-mesh}
  \end{subfigure}
  \hfill
  \begin{subfigure}[b]{0.48\textwidth}
    \includegraphics[width=\linewidth]{figures/occluder-splat}
    \caption{Photoreal twin.}
    \label{fig:occlusion-splat}
  \end{subfigure}
  \caption{The occlusion layer and the reconstruction it was derived from.
    (\subref{fig:occlusion-mesh}) The occluder: a decimated triangle mesh covering the
    buildings, trees, shrubs, and raised planting beds of the site. It is rendered here in
    grey for illustration only. In the deployed application it carries a depth-only
    material that writes to the depth buffer and outputs no color, so the layer is itself
    invisible while every virtual target and recall object behind it is hidden. Without
    it, targets on the far side of a building or a planting bed would remain visible
    through the geometry, and distance would carry no occlusion cue.
    (\subref{fig:occlusion-splat}) The same viewpoint in the Gaussian-splat
    reconstruction, from which the occluder was built. Because an occluder needs only the
    correct silhouette and not the correct surface, the mesh was reduced aggressively:
    geometry that occluded nothing was deleted, and buildings that did not reconstruct
    cleanly were replaced with boxes matching their footprint. Fidelity is spent on shape
    alone.}
  \Description{Two views of the same campus courtyard, side by side and at matching
    aspect ratio. The left image shows an untextured grey three-dimensional mesh on a dark
    background, its surface visibly faceted into coarse triangles: a long paved walkway
    runs from the lower left into the distance, with a raised curb along it, three rounded
    clumps of geometry standing in for shrubs in a planting bed, and a low ragged band of
    geometry across the background standing in for trees and buildings. Large flat areas
    are made of very few large triangles. The right image shows the same scene as a
    photorealistic reconstruction under a blue sky: a paved walkway and stone curb in the
    foreground, a green lawn, a bed of small shrubs and young trees, a red brick building
    with rows of windows along the right, a glass-fronted building at the left, and parked
    vehicles at the far right. The grey mesh on the left reproduces the silhouettes of the
    curb, shrubs, and buildings on the right, but none of their color or detail.}
  \label{fig:occlusion}
\end{figure*}

% -------------------------------------------------------------------------------------
% Single-column fallback, if the two panels crowd the layout. Same file names; stack the
% panels instead of placing them side by side. Delete the block above.
% -------------------------------------------------------------------------------------
%
% \begin{figure}[t]
%   \centering
%   \begin{subfigure}[b]{\linewidth}
%     \includegraphics[width=\linewidth]{figures/occluder-mesh}
%     \caption{Occlusion layer.}
%     \label{fig:occlusion-mesh}
%   \end{subfigure}
%   \\[2pt]
%   \begin{subfigure}[b]{\linewidth}
%     \includegraphics[width=\linewidth]{figures/occluder-splat}
%     \caption{Photoreal twin.}
%     \label{fig:occlusion-splat}
%   \end{subfigure}
%   \caption{The occlusion layer (\subref{fig:occlusion-mesh}) and the Gaussian-splat
%     reconstruction it was derived from (\subref{fig:occlusion-splat}). The occluder
%     carries a depth-only material in the deployed application: it writes depth and
%     outputs no color, so the layer is invisible while hiding every virtual object behind
%     it. Because only the silhouette matters, the mesh is decimated aggressively,
%     geometry that occluded nothing is deleted, and poorly reconstructed buildings are
%     replaced with footprint-matched boxes.}
%   \label{fig:occlusion}
% \end{figure}
